% TOP_PROBLEM: {"id": 2520, "title": "Kourovka Notebook Problem 21.11", "category_id": 17, "category": {"id": 17, "name": "group_theory", "display_name": "Group Theory", "description": "Problems about groups, group actions, representations, and related algebraic structures.", "slug": "group-theory", "order_index": 17, "created_at": "2026-07-31T15:26:25.671Z"}, "status": "open", "problem_number": "KOU-21.11", "set_id": 10, "statusQualification": "The two target classes, omitted from the catalogue statement field, are restored from its background and primary Problem21.11."}
% FIRST_POSTED: 2026-10-01
% ENTRY_KIND: statement_only
% UnsolvedMath ID 2520; source catalogue code KOU-21.11
% !TeX program = lualatex
% Definition and sources only; this file is not an LLM solution attempt.
\documentclass[11pt,a4paper]{article}
\usepackage[margin=25mm]{geometry}
\usepackage{fontspec}
\setmainfont{lmroman10-regular.otf}[BoldFont=lmroman10-bold.otf,
 ItalicFont=lmroman10-italic.otf,BoldItalicFont=lmroman10-bolditalic.otf]
\usepackage{amsmath,amssymb,mathtools}
\usepackage{unicode-math}
\setmathfont{latinmodern-math.otf}
\AtBeginDocument{\let\setminus\smallsetminus}
\usepackage{xurl,hyperref}
\hypersetup{colorlinks=true,urlcolor=blue}
\setcounter{secnumdepth}{0}
\setlength{\parindent}{0pt}
\setlength{\parskip}{5pt}
\setlength{\emergencystretch}{3em}
\begin{document}
\section*{Kourovka Notebook Problem 21.11}\label{2520}
{\small UnsolvedMath ID 2520 · KOU-21.11 · Group Theory · Kourovka Notebook - New Problems, Issue 21}
\subsection{Definitions and mathematical statement}
Write $\omega=\{0,1,2,\ldots\}$. A nonprincipal ultrafilter $\mathcal U$
on $\omega$ is a collection of subsets closed under finite intersections
and supersets, containing no finite set, such that exactly one of $A$
and $\omega\setminus A$ belongs to $\mathcal U$ for every $A\subseteq\omega$.
For a sequence of groups $(G_i)_{i\in\omega}$, its ultraproduct is
\[
 \prod_{i\in\omega}G_i/\mathcal U
 =\left(\prod_{i\in\omega}G_i\right)/
   \{(g_i):\{i:g_i=1\}\in\mathcal U\}.
\]
Multiplication is coordinatewise before passing to this normal-subgroup
quotient. A homomorphic image means the target of a surjective abstract
group homomorphism, with no continuity requirement.

Determine which groups in the following two classes can occur as such
images. Part (a) concerns infinite finitely generated groups of finite
exponent: a finite list generates the group, and some integer $n\ge1$
satisfies $g^n=1$ for every element $g$. Part (b) concerns
$\operatorname{FSym}(X)$ for an infinite set $X$, the group of all
bijections $X\to X$ having finite support $\{x:g(x)\ne x\}$.

For each proposed target, both the sequence of source groups and the
nonprincipal ultrafilter may be chosen. Countability refers to the
index set of the family; the individual groups $G_i$ need not be
countable or finite. The source asks whether some or all members of
each target class occur, so existence for a single target and occurrence
for every target are distinct parts of the requested determination.
\subsection{Short English statement}
Which infinite finitely generated finite-exponent groups and finitary symmetric groups are quotients of countably indexed nonprincipal ultraproducts?
\subsection{Sources}
\begin{itemize}
\item[S1] ULAM.AI and the UnsolvedMath Contributors, supplied problems.json, record 2520 (KOU-21.11), Kourovka Notebook - New Problems, Issue 21. Catalogue identity and source transcription; CC BY 4.0.
\url{https://huggingface.co/datasets/ulamai/UnsolvedMath}.
\item[S2] The Kourovka Notebook, 21st edition, new problems, Problem 21.11. Expanded formulation of the supplied catalogue statement.
\url{https://arxiv.org/pdf/1401.0300v47}.
\end{itemize}
\end{document}
